\documentclass[11pt,a4paper]{article}
\usepackage[margin=2.2cm]{geometry}
\usepackage{amsmath,amssymb}
\usepackage{booktabs}
\usepackage{longtable}
\usepackage{xcolor}
\usepackage[hidelinks]{hyperref}
\usepackage{fancyhdr}

\pagestyle{fancy}\fancyhf{}
\lhead{AUTHBC — Running Results Summary}\rhead{\thepage}
\newcommand{\ok}{\textcolor{green!55!black}{\textbf{PASS}}}
\newcommand{\bad}{\textcolor{red!80!black}{\textbf{FAIL}}}
\newcommand{\pend}{\textcolor{orange!85!black}{\textbf{pending}}}

\title{\textbf{AUTHBC — Running Results Summary}\\[2mm]
\large Co-optimizing encoding $\times$ authentication placement $\times$ signature scheme $\times$
batching\\ for blockchain-grade UAV telemetry over 802.11}
\author{Mohamed A. Farouk \\ \small (living document --- regenerated as each step completes)}
\date{\today}

\begin{document}
\maketitle
\thispagestyle{fancy}

\begin{abstract}
\noindent This document tracks every measured result in the AUTHBC project and is updated at each
step. All numbers trace to a frozen CSV in \texttt{results/} and are guarded by the reproduction gate
(\texttt{make verify-frozen}). The headline is established: the co-optimized configuration cuts
on-air authentication bytes by \textbf{75.0\,\%} versus the inline-CBOR baseline while holding
verifiability $V\ge0.95$ at loss $p=0.05$ --- criterion met. The hardware campaign has confirmed the
one predicted anomaly (the signature-scheme ordering flips on ARM) and has shown the energy model's
assumed CPU power to be $4.1\times$ too high.
\end{abstract}

\section{The problem}
Telemetry records are tiny (45--190\,B) but a signature is 64--96\,B, so authentication overhead
rivals the payload. Four knobs --- record \emph{encoding}, authentication \emph{placement} (A inline,
B self-batch, C relay-aggregate, D block), signature \emph{scheme}, and \emph{batch} size $b$ --- are
coupled and must be co-optimized. The falsifiable criterion: the co-optimized configuration must cut
on-air authentication bytes by $\ge 40\,\%$ versus inline-CBOR at $V\ge0.95$, $p=0.05$.

\section{Software layer (frozen, reproducible)}

\subsection{E1 --- overhead dominance (T1)}
$\varphi = g/(s+g)$ is the fraction of on-air bytes that is pure signature ($g=64$\,B).
\begin{center}\begin{tabular}{lrr}
\toprule encoding & mean record $s$ & $\varphi$ \\ \midrule
JSON & 191.1\,B & 25.1\,\% \\ CBOR & 66.25\,B & 49.1\,\% \\ MessagePack & 65.16\,B & 49.6\,\% \\
\textbf{delta} & \textbf{45.00\,B} & \textbf{58.7\,\%} \\ \bottomrule
\end{tabular}\end{center}
\textbf{Reading.} The better the compression, the \emph{worse} the relative overhead: delta is the
best encoder and leaves 59\,\% of the frame as signature. Compression alone cannot solve this --- the
motivation for the other three knobs.

\subsection{E2 --- batching cure (T2)}
Folding $b$ records under one signature gives $s+(g_a+H_f)/b$ per record; packing to the MTU bottoms
out at $A = M/(M-H_f-g_a) = 1.0745$ at $M=1500$. Measured: delta reaches $b_{\max}=31$ and attains
$A=1.0745$ exactly; other encodings within 5\,\% (integer-$b$ slack). Inline placement \emph{never}
amortizes its per-record signature --- only the header shrinks.

\subsection{E3 --- loss frontier (T3)}
Self-batch frames verify independently, so $V_B = 1-p$, \emph{flat in $b$}. A block signature
spanning $n$ frames needs all fragments: $V_D=(1-p)^n$. Measured at $p=0.10$, $b=40$:
$V_B=0.90$ versus $V_D=0.81=0.9^2$. This is the wall that stops over-aggregation.

\subsection{E4 --- scheme crossover (T4)}
BLS aggregates but its pairing verify is $10.7\times$ costlier than Ed25519. Which wins depends on the
powers only through $\kappa=P_r/P_c$, so the break-even $\kappa^\star$ is found without measuring
watts. \textbf{Result: Ed25519 wins all 80 grid points}, $\min\kappa^\star=3.21$. At the accepted
96\,B, BLS carries \emph{more} bytes than Ed25519's 64\,B on own traffic ($\kappa^\star=\infty$);
its aggregation pays off only for cross-signer relay traffic.

\subsection{E5 --- the headline (T5) \ok}
\begin{center}\begin{tabular}{llllrr}
\toprule config & enc & scheme & placement & auth B/rec & $V$ \\ \midrule
\textbf{optimized} & delta & \textbf{Ed25519} & B, $b{=}4$ & \textbf{26.00} & 0.95 \\
A+CBOR (baseline) & cbor & Ed25519 & A inline & 104.0 & 0.95 \\
A+JSON (naive) & json & Ed25519 & A inline & 104.0 & 0.95 \\
D-over-agg & cbor & Ed25519 & D, $b{=}40$ & 3.60 & \textbf{0.9025} \bad \\ \bottomrule
\end{tabular}\end{center}
\[ \text{cut}=\frac{104-26.0}{104}=\mathbf{75.0\,\%}\ \ge 40\,\%\quad\Rightarrow\quad \ok \]
Hand-checked: $45.0+104/4=71.0$\,B/record at $b{=}4$, freshness $4/20=200$\,ms.
\textbf{Freshness binds before the MTU does}: the byte-optimal batch is $b{=}31$ (96.8\,\%) but sits
at \textbf{1552\,ms} of staleness, $6.2\times$ over the $D_{\max}{=}250$\,ms bound docs/02 §7
requires the optimizer to enforce --- which it computed and discarded until audit \textbf{F10}.
Freshness is now a hard constraint and a fourth Pareto objective. The freshness-feasible batch obeys
$b\lesssim\Lambda D_{\max}$, independent of encoding and scheme. The resulting frontier
($b{=}1{\rightarrow}4$: $0/50/66.7/75\,\%$ cut at $50/100/150/200$\,ms) is the co-design result;
a single number is not. \textbf{The last row is the co-design argument in
miniature}: over-aggregation is byte-competitive (3.60\,B) yet \emph{fails} $V\ge0.95$ because
$0.9025=(1-p)^2$. The byte-optimal answer is not the correct answer.

\subsection{NS-3 channel validation}
\begin{center}\begin{tabular}{rrrr}
\toprule $N$ & unicast vs ACK-Bianchi & broadcast vs naive reduction & broadcast vs Ma \& Chen \\ \midrule
5 & $+0.61$\,\% & $-0.2$\,\% & $-0.33$\,\% \\
20 & $+0.64$\,\% & $+31.3$\,\% & $-0.71$\,\% \\
50 & $\mathbf{-2.85}$\,\% & $\mathbf{+1654.1}$\,\% & $\mathbf{+1.06}$\,\% \\ \bottomrule
\end{tabular}\end{center}
Unicast validates the Bianchi DCF model to $+0.6$--$-2.9\,\%$ (tightened from the previously
reported $\pm1.8$--$5.3\,\%$ by fixing F8, a $\sim$4.8\,\% goodput over-count in which the NS-3 sinks
outlived the sources by 0.5\,s on a 10\,s window, and by using exact OFDM symbol timing). The middle
column is our former in-house \emph{reduction} of the unicast model to broadcast ($\tau=2/(W{+}1)$);
it fails not because of capture --- measured at \textbf{0\,\%} --- but because it ignores the
\textbf{backoff counter Consecutive Freeze Process}: with no ACK the contention window never doubles,
so a station that has just transmitted may redraw backoff 0 and seize the medium one slot before any
deferring station, whose counter is necessarily $\ge1$. That is \textbf{Ma \& Chen's published
broadcast model} (IEEE Comm.\ Lett.\ 11(8):686--688, 2007; IEEE TVT 57(6):3757--3768, 2008), whose
abstract warns that unicast models ``cannot simply be reduced'' to broadcast --- exactly our error.
Their closed form reproduces our measurement to $\le0.36\,\%$ (F9). \textbf{No novelty is claimed for
the mechanism.}

\section{Hardware layer (Raspberry Pi 4B, measured)}

\subsection{Cross-platform determinism}
All \textbf{12/12} encoded-size metrics are byte-identical between x86 and ARM, as deterministic
integer-only encoders require.

\subsection{Timings and cycle normalisation}
\begin{center}\begin{tabular}{lrrrr}
\toprule op & x86 & RPi4 & wall ratio & \textbf{cycles ARM/x86} \\ \midrule
ed25519 sign & 29.1\,\textmu s & 88.1 & 3.03$\times$ & \textbf{1.16$\times$} \\
ed25519 verify & 95.0 & 259.5 & 2.73$\times$ & \textbf{1.05$\times$} \\
ecdsa sign & 26.2 & 135.5 & 5.17$\times$ & 1.98$\times$ \\
ecdsa verify & 78.5 & 326.9 & 4.16$\times$ & 1.60$\times$ \\
bls verify & 1016.5 & 8611.8 & 8.47$\times$ & 3.25$\times$ \\ \bottomrule
\end{tabular}\end{center}
Wall-clock ratio conflates clock speed with per-cycle efficiency. The clock floor is
$4.7/1.8=2.61\times$. Normalising shows Ed25519 needs only $1.05\times$ the cycles on ARM (equally
optimised on both ISAs) while BLS needs $3.25\times$ (heavily x86-tuned). \textbf{The docs/04 §1 ``5--15$\times$'' wall-clock anchor was therefore withdrawn} as
arithmetically unreachable and \textbf{restated per-cycle: expect 1--4$\times$ the CYCLES on ARM}
(clock- and platform-independent). All measured ops now sit inside that band.

\subsection{F6 --- the scheme ordering flips on ARM}
\begin{center}\begin{tabular}{lrrl}
\toprule platform & ECDSA verify & Ed25519 verify & winner \\ \midrule
x86 & 78.5\,\textmu s & 95.0\,\textmu s & ECDSA \\
\textbf{RPi4} & \textbf{326.9} & \textbf{259.5} & \textbf{Ed25519} \\ \bottomrule
\end{tabular}\end{center}
Three candidate causes were separated:
\begin{itemize}\itemsep2pt
\item \textbf{Library --- ruled out.} Byte-identical versions on both (\texttt{cryptography} 45.0.7).
\item \textbf{Our wrapper --- real but not the cause.} Our ECDSA path does ASN.1 DER conversion in
Python every op (Ed25519 does not): 1.05\,\textmu s on x86, \textbf{12.58\,\textmu s on ARM}
($\sim12\times$ worse). Flipping the result would have required $64\times$. Core-only on ARM: ECDSA
312.4 versus Ed25519 263.3 --- \textbf{Ed25519 still wins}.
\item \textbf{Hardware --- the operative cause}, but see \S\ref{sec:open}: the \emph{specific}
microarchitectural mechanism is \textbf{not} established.
\end{itemize}
\textbf{Consequence:} on the real deployment platform the optimum selects \textbf{Ed25519} (the D2
architectural default). Both are 64\,B, so the auth-byte headline is \textbf{unchanged} by the scheme flip.

\subsection{Measured per-op energy (RPi4, 60\,s windows, 5 reps)}
\begin{center}\begin{tabular}{lrrr}
\toprule op & energy/op (\textmu J) & 95\,\% CI & $\Delta P$ (W) \\ \midrule
ed25519 sign & \textbf{70.64} & --- & 0.773 \\
ed25519 verify & \textbf{187.47} & --- & 0.717 \\
ecdsa sign & 87.10 & --- & 0.634 \\
ecdsa verify & 201.04 & --- & 0.613 \\
bls sign & 1397.73 & --- & 0.513 \\
bls verify & 4555.00 & --- & 0.529 \\
delta encode & 36.92 & --- & 0.726 \\
cbor encode & 40.43 & --- & 0.753 \\ \bottomrule
\end{tabular}\end{center}
\textbf{Ed25519 is cheaper than ECDSA in energy as well as time on ARM} (70.6 vs 87.3\,\textmu J
sign; 186.8 vs 203.5 verify), consistent with the timing flip. Independent cross-check against
$\Delta P \times t_{op}$ (micro suite) agrees within 1--4.4\,\% for seven of eight ops;
all eight ops now agree within \textbf{3.7\,\%} (BLS to 0.0\,\% and $-0.1$\,\%). The earlier
14.4\,\% \texttt{bls:verify} discrepancy was traced to warmup running \emph{inside} the sync window
and was fixed at source and \textbf{re-measured} (superseded data retained under
\texttt{results/hw/energy/superseded/}). Measured $p_{cpu}$ median \textbf{0.634\,W}
(range 0.513--0.773).

\subsection{Reproducibility and energy}
Two boards, independent runs: worst-case divergence \textbf{0.5\,\%} across all 16 timed ops. The
energy chain validated end-to-end (12 windows $\to$ 12 segments) with two independent cross-checks
(loop throughput versus micro suite within 5--6\,\%; single-core $\Delta P$ 0.728\,W versus 4-core
stress 0.690\,W/core). \textbf{Measured $p_{cpu}\approx0.73$\,W against the nominal 3.0\,W used in
E5 --- $4.1\times$ too high.} The auth-byte headline is power-free and unaffected.

\section{Open items and inconsistencies}\label{sec:open}
\begin{longtable}{p{3.2cm}p{11.5cm}}
\toprule \textbf{item} & \textbf{status} \\ \midrule
F6 mechanism & \textbf{Claim retracted.} Masking ADX/BMI2 on x86 via
\texttt{OPENSSL\_ia32cap} did \emph{not} slow ECDSA (61.1 $\to$ 56.8\,\textmu s; ECDSA still won), so
the ADX explanation is \textbf{unsupported}. What remains established: the ordering is ISA-dependent
and Ed25519 is markedly more portable (1.05$\times$ versus 1.60$\times$ cycles). The specific
mechanism is an open question. \\ \addlinespace
NS-3 broadcast & \textbf{CLOSED (F9).} PHY instrumentation shows capture contributes
\textbf{0\,\%}: all 1149 successful decodes at the monitored receiver came from single-transmitter
busy periods, and an equal-power channel is byte-identical to the default. The cause is the
Consecutive Freeze Process, and the corrected model is \textbf{published} (Ma \& Chen 2007/2008):
their closed form matches our NS-3 measurement to $\le0.36\,\%$ at every $N$. Our contribution here
is validation at $W_0{=}16$ / 802.11a (they tested $W_0{=}32,128$ at 1\,Mb/s) and a direct PHY-trace
measurement of the mechanism. \\ \addlinespace
$\kappa$ plausible band & E4 assumes $\kappa=P_r/P_c\lesssim0.5$ from ``radio below CPU power''. With
$p_{cpu}$ measured at 0.73\,W (not 3.0\,W), $\kappa$ could approach or exceed 1. The E4 verdict is
unchanged ($\min\kappa^\star=3.21$) but the stated band must be re-derived from measured powers. \\
\addlinespace
$p_{radio}$ & \pend{} --- needs the 2-node link run; both boards now healthy. \\
E4/E5 re-run & \pend{} --- deliberately deferred until both powers are measured, so the energy column
stays internally consistent rather than half-measured. \\
5--15$\times$ anchor & Needs restating per-cycle (decision required). \\ \bottomrule
\end{longtable}

\section*{Provenance}
Every number above traces to a frozen CSV under \texttt{results/}; derived artefacts are regenerated
byte-identically by \texttt{make verify-frozen}, which fails loudly on any drift. Full audit:
\texttt{docs/audits/p7.md}; decisions and their limitations: \texttt{docs/DECISIONS.md}.

\end{document}
